Economics has long relied on dynamic programming to formalize and solve problems of sequential decision-making. Yet, for decades, the complexity of our models has been dictated less by the economic phenomena we wish to study and more by the computational limits of the curse of dimensionality. Reinforcement learning offers one pragmatic path forward. By replacing exhaustive state-space enumeration with sampling and function approximation, RL trades the absolute convergence guarantees of classical dynamic programming for the capacity to operate in high-dimensional, continuous, and complex strategic environments.

As this survey has sought to illustrate, RL is not a departure from dynamic programming, but a scalable and natural extension of it. It does come with its own issues. Deep reinforcement learning, in particular, remains brittle, sensitive to hyperparameters, and vulnerable to the structural instabilities of the deadly triad. Furthermore, its success in structural estimation and policy design relies heavily on computational tricks and simplifying assumptions, the availability of highly accurate simulators, and clearly defined transition or reward functions.

However, the intersection of these two fields is deeply complementary. While RL equips economics with new tools, economics provides RL with much-needed structure. As demonstrated in applications ranging from dynamic pricing to causal inference, incorporating fundamental economic concepts, such as revealed preference, parametric assumptions, or causal identification, can tame RL’s sample complexity and prevent algorithms from converging on systematically biased policies.

Looking ahead, we should view reinforcement learning simply as another tool in the applied economist’s toolkit. As these algorithms mature and their theoretical foundations solidify, the novelty of using "AI" to solve economic models will likely fade. However, one can expect interesting cross fertilization between these computational fields, that gradually expands the frontier of the models we can solve, the mechanisms we can design, and the behaviors we can understand.
